
\documentclass{beamer}
\usetheme{PaloAlto}
\usecolortheme{beaver}
\logo{\includegraphics[height=1.5cm]{vut_logo.eps}}

\usepackage[czech]{babel}
\usepackage[utf8]{inputenc}
\usepackage[T1]{fontenc}
\usepackage{times}
\usepackage{graphics}
\usepackage{picture}
\usepackage{marvosym} % smajlici
\definecolor{dkbrown}{RGB}{122,0,0}
\setbeamercolor{palette sidebar quaternary}{use=structure,fg=dkbrown}
\setbeamercolor{palette sidebar secondary}{use=structure,fg=dkbrown}
\setbeamercolor{item projected}{bg=dkbrown}

\title{Pojednání o~tečce}
\subtitle{aneb její nepostradatelnost}
\author[]{Marek Šipoš}
\institute{Vysoké učení technické v~Brně\\Fakulta informačních technologií}
\date{\today}

% --- ZVÝRAZNĚNÝ ÚSEK TEXTU ---
% \alert{highlighted}

% --- BOXÍKY ---
% \begin{block}{Remark}   \end{block}
% \begin{alertblock}{Important theorem}   \end{alertblock}
% \begin{examples}   \end{examples}

% --- SLOUPCE ---
% \begin{columns}   \column{0.5\textwidth}   \end{columns}

% --- ANIMACE ---
% \item<1->
% \pause

\begin{document}
    %%%%%%%%%%%%%%%%%%%%%%%%%%%%%%%%%%%%%%%%%%%%%%%%%
    \section[Úvod]{}
    \frame{\titlepage}
    %%%%%%%%%%%%%%%%%%%%%%%%%%%%%%%%%%%%%%%%%%%%%%%%%
    \section[Tečka]{}
    \begin{frame}
        \frametitle{Tečka}
        {\fontsize{1cm}{1cm}\selectfont $\ldotp$}\hspace{5mm}{\large Tečka,} {bodka,} {\small period,} {\footnotesize piste,} {\scriptsize pont,} {\tiny kropka}\\
        \vspace{1cm}
        \begin{itemize}
            \item<1->Různé názvy, stejný význam
            \item<2->Ale jaký je to vlastně význam?
            \item<3->Význam tečky se vyvíjel napříč historií
        \end{itemize}
        
    \end{frame}
    %%%%%%%%%%%%%%%%%%%%%%%%%%%%%%%%%%%%%%%%%%%%%%%%%
    \section[Odkud pochází]{}
    \begin{frame}
        \frametitle{Odkud tečka pochází?}
        Kdy se tečka vyskytla poprvé?\\
        \vspace{1cm}
        \pause
        \alert{První známý výskyt:}\\
        \begin{itemize}
            \item<3->3. století př. n. l.
            \item<4->Vzešla z~řecké interpunkce
            \item<5->Její význam určovalo její umístění
        \end{itemize}
        
    \end{frame}
    \begin{frame}
        \frametitle{Odkud tečka pochází?}
        Poziční sémantika tečky ve starověkém Řecku:\\
        \vspace{5mm}
        \begin{columns}
            {\fontsize{7mm}{1cm}\selectfont $^\ldotp$}
            \column{0.9\textwidth}
            \uv{Vysoká tečka}\,--\,zakončovala myšlenku
        \end{columns}\vspace{2mm}
        \begin{columns}
            {\fontsize{5mm}{1cm}\selectfont $\ldotp$}
            \column{0.9\textwidth}
            \uv{Střední tečka}\,--\,delší nádech
        \end{columns}\vspace{2mm}
        \begin{columns}
            {\fontsize{7mm}{1cm}\selectfont $_\ldotp$}
            \column{0.9\textwidth}
            \uv{Nízká tečka}\,--\,kratší nádech
        \end{columns}
    \end{frame}
    %%%%%%%%%%%%%%%%%%%%%%%%%%%%%%%%%%%%%%%%%%%%%%%%%
    \section[Různé druhy]{}
    \begin{frame}
        \frametitle{Různé druhy \uv{teček}}
        Druhy teček na základě grafické reprezentace\\
        \vspace{1cm}
        \pause
        \alert{O~tečce můžeme mluvit jako o:}\\
        \begin{itemize}
            \item<3->Běžné tečce tak, jak ji známe
            \item<4->Puntíku (např. typ odrážky v~seznamu)
            \item<5->Bodě (čistě grafická reprezentace v~prostoru)
        \end{itemize}
    \end{frame}
    %%%%%%%%%%%%%%%%%%%%%%%%%%%%%%%%%%%%%%%%%%%%%%%%%
    \section[Základ abeced]{}
    \begin{frame}
        \frametitle{Tečka jako základ abeced}
        Tečka je nepostradatelným prvkem některých abeced
    \end{frame}
    \begin{frame}
        \frametitle{Tečka jako základ abeced}
        \alert{Braillovo písmo}\\
        \begin{itemize}
            \item Počátek 19. století
            \item Body vyražené do materiálu (ve skupinách po šesti)
            \item Bod je zvolen z~praktických důvodů
            \begin{itemize}
                \item Jednoduchá realizace
                \item Snadné vnímání hmatem
            \end{itemize}
        \end{itemize}
    \end{frame}
    \begin{frame}
        \frametitle{Tečka jako základ abeced}
        \alert{Morseova abeceda}\\
        \begin{itemize}
            \item Konec 19. století
            \item Tečka jako krátký signál
            \begin{itemize}
                \item podobnost s~tečkou jako krátkým nádechem v~Řecku?
            \end{itemize}
            \item Tečka samotná značí písmeno 'E'
            \item Může to být i~zkrácený kód pro číslici '5'
            \item Tečka samotná se však zapisuje \texttt{.-.-.-}
        \end{itemize}
    \end{frame}
    \begin{frame}
        \frametitle{Tečka jako základ abeced}
        \alert{Brainf***}\\
        \begin{itemize}
            \item Vznik 1993
            \item Minimalistický ezoterický programovací jazyk
            \item Tečka jako příkaz pro výpis buňky na standardní výstup
        \end{itemize}
    \end{frame}
    %%%%%%%%%%%%%%%%%%%%%%%%%%%%%%%%%%%%%%%%%%%%%%%%%
    \section[Interpunkce]{}
    \begin{frame}
        \frametitle{Interpunkční znaménko}
        Tečka se hojně využívá jako interpunkční znaménko s~různými významy:\\
        \vspace{3mm}
        \begin{itemize}
            \item<2-> Zakončení vět
            \begin{itemize}
                \item<3-> Oznamovacích vět
                \item<4-> Ale ani znaky '?' '!' se bez tečky neobejdou \Smiley{}
            \end{itemize}
            \item<5-> Iniciály ({\em George W. Bush})
            \item<6-> Zkratky, tituly, atd.
            \item<7-> Výpustek (\dots)
        \end{itemize}
    \end{frame}
    %%%%%%%%%%%%%%%%%%%%%%%%%%%%%%%%%%%%%%%%%%%%%%%%%
    \section[Technika]{}
    \begin{frame}
        \frametitle{Využití v~technice}
        Tečka je nepostradatelná v~různých technických oblastech
    \end{frame}
    \begin{frame}
        \frametitle{Využití v~technice}
        \begin{columns}
            \column{0.5\textwidth}
            \alert{Matematika}\\
            \begin{itemize}
                \item<2-> Desetinná tečka\\\texttt{3.14159\dots}
                \item<3-> Oddělovač\\\texttt{42.000}
                \item<4-> Středová tečka\\\texttt{3$\cdot$3=9}\\\texttt{$\vec{u}\cdot\vec{v}$}
            \end{itemize}
            \column{0.5\textwidth}
            \pause[5]
            \alert{Informatika}\\
            \begin{itemize}
                \item<6-> Pro přístup\\\texttt{myStruct.something}
                \item<7-> Oddělovač\\\texttt{www.wikipedia.org}\\\texttt{127.0.0.1}
                \item<8-> Nadř. adresář\\\texttt{cd $..$}
            \end{itemize}
        \end{columns}
        \pause[9]
        \vspace{4mm}
        To byl jen zlomek z~případů využití tečky!
    \end{frame}
    %%%%%%%%%%%%%%%%%%%%%%%%%%%%%%%%%%%%%%%%%%%%%%%%%
    \section[Konec]{}
    \begin{frame}
        \begin{center}
            \frametitle{Konec}
            \alert{A~to je konec\dots}\\
            \vspace{1cm}
            Já i~tečka děkujeme za pozornost \Smiley{}
        \end{center}
    \end{frame}
    %%%%%%%%%%%%%%%%%%%%%%%%%%%%%%%%%%%%%%%%%%%%%%%%%
\end{document}